\scanpage{134}
% archive sheet number: 133

$\simeq(\det H)\otimes\det(H')$. Cette \deleted{\unsure{loi}} fonction
quadratique est \uline{universelle}, pour
les fonctions quadratiques

\diagram{\begin{tikzcd}[ampersand replacement=\&]
P\arrow[r,"q"] \& Q
\end{tikzcd}}

telles que $q(x\otimes y)=0\qquad\forall x,y\in P$.

Ceci nous montre \deleted{que, si} (revenant
à \deleted{$\unsure{Q_P|\underline Q_P}$} la situation \unsure{précédente|précédent}) que
si $\underline Q_P$ n'est pas vide, c'est un torseur
sous \deleted{$\unsure{\underline L_B|\underline L_P}$} $(L\otimes L')^\vee$. Il faut vérifier
que $\alpha,\delta$ sont compatibles avec
structure torseur -- moyennant quoi
on voit que c'est des iso.

Pour vérifier que $\underline Q_P$ n'est pas
vide, on \deleted{prend sur $P$ la structure
\deleted{\unsure{tensorielle}} d'alg. produit tensoriel
de celles de $H,H'$, \deleted{et la fon} trouve une Alg. de H.C., dont
la fonction $\det$ fait l'affaire.}

on peut utiliser la construction
des algèbres $M(b,c,b',c',\underline t)$. Cela
me fait songer à un sixième
ensemble

